\section{Материал, не вошедший в лекции}
\subsection{Неполнота $RL_p$}
\begin{definition}
    Пусть $RL_p([a,b])$ — линейное пространство функций, $p$-ая степень модуля которых интегрируема по Риману.
\end{definition} 

\begin{theorem}
    Пространство $R L_p([a, b])$ не полно.
\end{theorem} 
\begin{proof}
    Без ограничения общности, $[a, b]=[0,1]$.
    Пусть $\epsilon \in\left(0, \frac{1}{2}\right)$.
    Перенумеруем рациональные точки отрезка $[0,1]: \Q \cap[0,1]=\left\{r_k\right\}$.
    $$
    G_n:=\bigcup_{k=1}^n\left(r_k-\frac{\varepsilon}{2^{k+2}}, r_k+\frac{\varepsilon}{2^{k+2}}\right) \cap[0,1],
    $$
    $$G:=\bigcup_{n=1}^{\infty} G_n.$$

    Тогда $\chi_{G_n}$ интегрируема по Риману по критерию Лебега, потому что как характеристическая функция объединения конечного набора интервалов, пересечённых с отрезком, она обладает конечным числом разрывов. Таким образом,
    \begin{equation*}
    \left\{\chi_{G_n}\right\} \subset R L_p([0,1]).
    \end{equation*}

    Докажем, что $\chi_G$ имеет множество точек разрыва положительной меры Лебега. Для этого рассмотрим $F=[0,1] \backslash G$. Тогда $\chi_G(F)=0$, но в то же время $\chi_G(\Q \cap[0,1])=1$, а так как $\Q$ плотно в $\R$, во всех точках $F$ функция $\chi_G$ разрывна. При этом по счётной полуаддитивности меры Лебега $\mathcal{L}^n(G) \leqslant \frac{1}{2}$, а значит $\mathcal{L}^n(F) \geqslant \frac{1}{2}>0$. Итак, $\chi_G$ разрывна на множестве $F$ положительной меры, итого
    \begin{equation*}
    \chi_G \notin R L_p([0,1]).
    \end{equation*}

    Введём обозначения
    $$
    E_k:=\left(r_k-\frac{\varepsilon}{2^{k+2}}, r_k+\frac{\varepsilon}{2^{k+2}}\right) \cap[0,1],
    $$
    $$G_m^n:=\bigcup_{k=n}^m E_k$$

    (в новых обозначениях $G_n=G_n^1$) и покажем фундаментальность последовательности $\left\{\chi_{G_n}\right\}$:
    \begin{multline*}
    \int_0^1\left|\chi_{G_m}(x)-\chi_{G_n}(x)\right| d x
    =
    (\text {так как } G_n \subseteq G_m)
    =\\=
    \int_0^1 \chi_{G_m \backslash G_n}(x) d x \leqslant \int_0^1 \chi_{G_m^n}(x) d x \leqslant \\
    \leqslant \sum_{k=n+1}^m \int_0^1 \chi_{E_k}(x) d x=\sum_{k=n+1}^m \frac{\varepsilon}{2^k} \rightarrow 0, \text { при } n, m \rightarrow \infty
    \end{multline*} 

    Итого, последовательность $\left\{\chi_{G_p}\right\}$ фундаментальна в $R L_p([0,1])$.

    Как известно, $R L_p$ — подпространство $L_p$. Тогда если у последовательности из $R L_p$ есть предел в $R L_p$, то он такой же и при рассмотрении её как последовательности в $L_p$ с пределом в $L_p$. В нашем случае, если бы у $\left\{\chi_{G_n}\right\}$ был предел в $R L_p([0,1])$, то у неё был бы такой же предел в $L_p([0,1])$. Но мы знаем, что её предел в $L_p([0,1])$ — это $\chi_G$, которая не лежит в $R L_p([0,1])$. Значит, у $\left\{\chi_{G_n}\right\}$ нет предела в $R L_p([0,1])$, однако это фундаментальная в $R L_p([0,1])$ последовательность. Значит пространство $R L_p([0,1])$ не полно.
\end{proof} 


\subsection{Функции ограниченной вариации}
\begin{definition}
    Пусть $T$ — разбиение отрезка $[a, b]$, т.е.
    $$T=\left\{x_i\right\}_{i=0}^{N_T}, \quad N_T \in \mathbb{N},$$
    $$a=x_0<x_1<\cdots<x_{N_T}=b.$$
    
    Пусть $f:[a, b] \rightarrow \mathbb{R}$. Тогда $V_T(f)$ — \textit{вариация функции $f$ по разбиению $T$}:
    $$V_T(f) :=\sum_{k=0}^{N_T-1}\left|f\left(x_{k+1}\right)-f\left(x_k\right)\right|,$$
    $$V_a^b(f) :=\sup _{T-\text { разбиение }[a, b]} V_T(f).$$
\end{definition} 

\begin{definition}
    $f$ называется \textit{функцией ограниченной вариации на $[a, b]$}, если
    $$V_a^b(f)<+\infty$$

    Обозначается $f \in B V([a, b])$.
\end{definition} 

\begin{theorem}
    $BV ([a,b])$ — линейное пространство.
\end{theorem} 
\begin{proof}
    Покажем, что $f_1, f_2 \in B V([a, b]) \Longrightarrow \alpha f_1+\beta f_2 \in B V([a, b])$. Пусть $T$ — произвольное разбиение $[a, b]$. Тогда по неравенству треугольника:
    $$V_T\left(\alpha f_1+\beta f_2\right) \leqslant|\alpha| V_T\left(f_1\right)+|\beta| V_T\left(f_2\right).$$

    Взяв супремум по всем разбиениям $[a,b]$, получим желаемое.
\end{proof} 

\begin{lemma}
    Если $\forall f:[a, b] \rightarrow \R$ монотонна на $[a, b]$, то $f \in B V([a, b])$ и её $V_a^b(f)=|f(b)-f(a)|$.
\end{lemma} 
\begin{proof}
    Все $V_T(f):=\sum_{k=0}^{N_T-1}\left|f\left(x_{k+1}\right)-f\left(x_k\right)\right|$ одного знака, то есть модуль раскрывается одинаково, а тогда соседние слагаемые друг друга уничтожают.
\end{proof}

\begin{lemma}
    Пусть $-\infty<a<c<b<+\infty$. Тогда
    $$
    f \in B V([a, b]) \Longleftrightarrow\left\{\begin{array}{l}
    f \in B V([a, c]), \\
    f \in B V([c, b]).
    \end{array}\right.
    $$
    B случае, если $f \in B V([a, b])$, тогда
    $$V_a^b(f)=V_a^c(f)+V_c^b(f).$$
\end{lemma} 
\begin{proof}\tab
    \begin{itemize}
        \item[$\underline{\Longrightarrow}$] Пусть $f \in B V([a, b])$. Обозначим:
        \begin{enumerate}
            \item $T_1$ — произвольное разбиение $[a, c]$,
            \item $T_2$ — произвольное разбиение $[c, b]$,
            \item $T=T_1 \cup T_2$ — разбиение отрезка $[a, b]$.
        \end{enumerate}
        Тогда имеем, что $V_{T_1}(f)+V_{T_2}(f) \leqslant V_T(f) \leqslant V_a^b(f)$.
        Взяв супремум сначала по $T_1$, а потом по $T_2$, получим
        \[V_a^c(f)+V_c^b(f) \leqslant V_a^b(f).\]
    \item[$\underline{\Longleftarrow}$]
        Пусть 
        $$\left\{\begin{array}{l}f \in B V([a, c]), \\ f \in B V([c, b]).\end{array}\right.$$
        Пусть $T=\left\{x_i\right\}_{i=0}^N$ — произвольное разбиение отрезка $[a, b]$.
        Если $c=x_{i^*}$ при некотором $i^*$, то это простой случай, так как тогда можно $\left\{x_j\right\}_{j=0}^{i^*}$ выбрать в качестве $T_1$, а $\left\{x_j\right\}_{j=i}^N$ выбрать в качестве $T_2$. И тогда очевидным образом $V_T(f)=V_{T_1}(f)+V_{T_2}(f) \leqslant V_a^c(f)+V_c^b(f)<+\infty$, а взяв супремум по всем $T$, получим $V_a^b(f) \leqslant V_a^c(f)+V_c^b(f)<+\infty$.

        Теперь рассмотрим случай, когда ни при каком $i$ $x_i$ не равно $c$. Тогда $c \in \left(x_i, x_i+1\right)$ при некотором $i$:
        \begin{multline*}
            V_T(f) =\sum_{k=0}^{N-1}\left|f\left(x_{k+1}\right)-f\left(x_k\right)\right|
            =\\=
            \sum_{k=0}^{i-1}\left|f\left(x_{k+1}\right)-f\left(x_k\right)\right|
            +
            \left|f\left(x_i\right)-f\left(x_{i+1}\right)\right|
            +
            \sum_{k=i+1}^{N-1}\left|f\left(x_{k+1}\right)-f\left(x_k\right)\right| \leqslant \\
            \leqslant \sum_{k=0}^{i-1}\left|f\left(x_{k+1}\right)-f\left(x_k\right)\right|
            +\\+
            \left|f\left(x_i\right)-f(c)\right|+\left|f(c)-f\left(x_{i+1}\right)\right|+\sum_{k=i+1}^{N-1}\left|f\left(x_{k+1}\right)-f\left(x_k\right)\right| \leqslant (*)
        \end{multline*} 
        Обозначим разбиения:
        $$
        \begin{aligned}
             T_1=\left\{x_0, x_1, \ldots, x_i, c\right\}, \\
            T_2=\left\{c, x_{i+1}, \ldots, x_N\right\}.
        \end{aligned}
        $$
        Тогда полученный ранее результат можно оценить как
        $$(*) \leqslant V_{T_1}(f)+V_{T_2}(f) \leqslant V_a^c(f)+V_c^b(f).$$

        Итого $V_T(f) \leqslant V_a^c(f)+V_c^b(f)$. Взяв супремум по всем $T$, получим:
        $$
        \begin{gathered}
        \sup _T V_T(f) \leqslant V_a^c(f)+V_c^b(f), \\
        V_a^b(f) \leqslant V_a^c(f)+V_c^b(f).
        \end{gathered}
        $$
        Если оба слагаемых в правой части конечны, то и $V_a^b$ конечна. Итого из первого и второго пункта
        $$
        V_a^b(f)=V_a^c(f)+V_c^b(f).
        $$
    \end{itemize}
\end{proof} 

\begin{theorem}
    Пусть $f \in B V([a, b])$. Тогда функция $g(x):=V_a^x$ монотонно не убывает на $[a, b]$. 
\end{theorem}
\begin{proof}
    Пусть $x_2>x_1$. Тогда применим только что доказанную лемму, выбрав $a= a, c=x_1, b=x_2$:
    $$
    \begin{gathered}
    V_a^{x_2}(f)=V_a^{x_1}(f)+V_{x_1}^{x_2}(f), \\
    V_a^{x_2}(f)-V_a^{x_1}(f)=V_{x_1}^{x_2}(f) \geqslant 0, \\
    g\left(x_2\right)-g\left(x_1\right) \geqslant 0, \\
    g\left(x_2\right) \geqslant g\left(x_1\right).
    \end{gathered}
    $$
\end{proof} 

\begin{theorem}
    Пусть $f \in B V([a, b])$. Тогда существуют $f_1$ и $f_2$ монотонно неубывающие на $[a, b]$ такие, что $f=f_1-f_2$.
\end{theorem} 
\begin{proof}
    Определим $f_1(x):=V_a^x(f) \quad \forall x \in[a, b]$. По только что доказанной теореме это монотонно неубывающая функция.
    Докажем, что функция $f_2(x)=f_1(x)-f(x)$ монотонно не убывает. Пусть

    $$a \leqslant x \leqslant y \leqslant b.$$ 
    \begin{multline*}
        f_2(y)-f_2(x)=\left[f_1(y)-f(y)\right]-\left[f_1(x)-f(x)\right] =\\=\left[f_1(y)-f_1(x)\right]-[f(y)-f(x)] \stackrel{(1)}{=} V_x^y(f)-[f(y)-f(x)].
    \end{multline*} 

    (1) — в силу аддитивности вариации по отрезкам.

    Заметим, что $V_x^y(f)=\sup_T V_T(f) \geqslant V_{\{x, y\}}(f)=|f(y)-f(x)|$. Тогда предыдущее выражение не меньше 0, а значит $f_2$ не убывает.
\end{proof} 
\begin{corollary}
    $\forall f \in B V([a, b])$ имеет не более чем счётное множество точек разрыва, при том все они 1-го рода. (Из 1 семестра помним, что монотонная функция имеет не более чем счетное число точек разрыва, при том все они первого рода).
\end{corollary}


\subsection{Абсолютно непрерывные функции}
\begin{definition}
    Функция $f:[a, b] \rightarrow \mathbb{R}$ \textit{называется абсолютно непрерывной на $[a, b]$}, если
$$
\forall \epsilon \ \exists \delta(\epsilon)>0: \forall \text { дизъюнктивной системы }\left\{\left(a_k, b_k\right)\right\}_{k=1}^N: \sum_{k=1}^N\left|a_k-b_k\right|<\delta(\epsilon)
$$
верно
$$\sum_{k=1}^n\left|f\left(b_k\right)-f\left(a_k\right)\right|<\varepsilon.$$
\end{definition} 

$AC([a, b])$ — множество всех абсолютно непрерывных на $[a, b]$ функций.

\begin{remark}
    $f \in A C([a, b])$ является непрерывной на $[a, b]$, поскольку из абсолютной непрерывности следует равномерная непрерывность. Обратное неверно.
\end{remark}

\noindent\textbf{Контрпример.}
$$
f(x)=\left\{\begin{array}{ll}
x \sin \frac{1}{x}, & x \neq 0 \\
0, & x=0.
\end{array}\right.
$$

$f \in C([0,1])$, но $f \notin B V([0,1])$, а значит по теореме, которая будет доказана ниже, $f \notin A C([0,1])$.

\begin{theorem}
    Если $f \in A C([a, b])$, то $f \in B V([a, b])$.
\end{theorem} 
\begin{proof}
    Запишем определение абсолютной непрерывности при $\epsilon=1$. Тогда $\exists \delta=\delta(1)> 0$: для любого конечного попарно непересекающегося набора интервалов $\left\{\left(a_k, b_k\right)\right\}_{k=1}^N: \sum_{k=1}^N\left|b_k-a_k\right|<\delta$ верно
    $$
    \sum_{k=1}^N\left|f\left(b_k\right)-f\left(a_k\right)\right|<1 .
    $$

    Теперь разобьём отрезок $[a, b]$. Пусть $\N \ni M:=\left[\frac{|b-a|}{\delta}\right]+1$.
    $$
    \begin{aligned}
    & x_0=a \\
    & x_1=a+\frac{b-a}{M} \\
    & \vdots \\
    & x_M=a+b-a=b
    \end{aligned}
    $$

    То есть поделили отрезок $[a, b]$ на одинаковые куски.
    Так как $x_{i+1}-x_i<\delta$ (специально для выполнения этого было взято достаточно большое $M$), то любое разбиение отрезка $\left[x_i, x_{i+1}\right]$ образует естественным образом конечный набор дизъюнктных интервалов суммарной длины меньше $\delta$, а значит по абсолютной непрерывности $V_{x_i}^{x_i+1}(f)<1 \quad \forall i$ Тогда в силу аддитивности вариации
    $$
    V_a^b(f)=\sum_{i=0}^{M-1} V_{x_i}^{x_{i+1}}(f)<\sum_{i=0}^{M-1} 1=M<+\infty.
    $$
\end{proof} 


\subsection{$N$-свойство Лузина}
\begin{definition}
    Пусть $f: \R^n \to \R$. Будем говорить, что функция обладает \textit{$N$-свойством Лузина}, если для любого измеримого множества $E \subset \R^n$ верно 
    $$\mathcal{L}^n(E) = 0 \implies \mathcal{L}^1(f(E))=0.$$
\end{definition} 

\begin{theorem}
    Пусть $f \in AC([a,b])$. Тогда оно обладает $N$-свойством Лузина.
\end{theorem} 
\begin{proof}
    Пусть $E \in[a, b], \mathcal{L}^1(E)=0$. В силу регулярности меры Лебега $\forall \delta>0$ существует открытое множество $E_\delta \supset E$ и $\mathcal{L}^1\left(E_\delta\right)<\delta$.
    Воспользуемся фактом:

    \noindent\textbf{Факт.} Любое открытое множество $G$ на числовой прямой может быть представлено в виде не более, чем счётного объединения попарно непересекающихся интервалов.

    Тогда $\forall \delta>0$ $E_\delta=\bigsqcup_{k=1}^{\infty}\left(a_k(\delta), b_k(\delta)\right) \supset E$. Получаем, что 
    $$\sum_{k=1}^{\infty}\left|b_k(\delta)-a_k(\delta)\right|<\delta.$$
    Заметим, что 
    $$f\left(\left[a_k(\delta), b_k(\delta)\right]\right)=\left[c_k(\delta), d_k(\delta)\right]$$ 
    — образ отрезка есть отрезок. Кроме того,
    $$
    f(E) \subset f\left(E_\delta\right)=\bigcup_{k=1}^{\infty} f\left(\left[a_k(\delta), b_k(\delta)\right]\right)=\bigcup_{k=1}^{\infty}\left[c_k(\delta), d_k(\delta)\right]
    $$

    Выберем $x_k(\delta) \in\left[a_k(\delta), b_k(\delta)\right]$ и $y_k(\delta) \in\left[a_k(\delta), b_k(\delta)\right]$ такие, что $f\left(x_k(\delta)\right)=c_k(\delta)$ и $f\left(y_k(\delta)\right)=d_k(\delta)$. 
    Но тогда интервалы ($x_k(\delta), y_k(\delta)$) попарно не пересекаются. Но так как $f \in A C([a, b])$, то $\forall \epsilon>0$ $\exists \delta(\epsilon)$ такой, что для любой системы из попарно непересекающися интервалов, суммарная разница длин которых не более $\delta(\epsilon)$, выполнено, что сумма модулей разности значений меньше $\epsilon$. $\forall \epsilon>0$ существует дизъюнктный набор интервалов $\left\{\left(x_k(\delta(\epsilon)), y_k(\delta(\epsilon))\right)\right\}_{k=1}^{\infty}$ и $\sum_{k=1}^{\infty}\left|x_k(\delta(\epsilon))-y_k(\delta(\epsilon))\right|<\delta(\epsilon)$. Но в силу абсолютной непрерывности $f$:
    $$
    \sum_{k=1}^{\infty}\left|f\left(x_k(\delta(\epsilon))\right)-f\left(y_k(\delta(\epsilon))\right)\right|=\sum_{k=1}^{\infty}\left|c_k(\delta(\epsilon))-d_k(\delta(\epsilon))\right|<\epsilon .
    $$

    Получаем, что $\mathcal{L}_*^1(f(E))<\epsilon$ — внешняя мера. Но так как мера Лебега полна, а $\epsilon$ выбрано произвольно, то $f(E)$ измеримо и $\mathcal{L}^1(f(E))=0$.
\end{proof} 


\subsection{Теорема Банаха-Зарецкого}
\begin{theorem}[Банах, Зарецкий]
    $$
    f \in A C([a, b]) \Longleftrightarrow\left\{\begin{array}{l}
    f \in B V([a, b]) \\
    f \in C([a, b]) \\
    \text { f обладает N-свойством Лузина }
    \end{array}\right.
    $$
\end{theorem} 
\begin{proof}\tab
    \begin{itemize}
        \item[$\underline{\Longrightarrow}$] Воспользуемся результатами предыдущих теорем.
        \item[$\underline{\Longleftarrow}$] Без доказательства.
    \end{itemize}
    
\end{proof} 

Приведём пример функции из $BV([a,b])$ и $C([a,b])$, но при этом не обладающей $N$-свойством Лузина.

\begin{example}[Канторова лестница]
    Построим канторово множество. Будем строить итеративно: на каждой итерации будем брать один из отрезков, полученных на предыдущей итерации, делить его на три равных части и «выбрасывать» среднюю. На $k$-ом шаге у нас будет $2^k$ равных отрезков длины $\left(\frac{2}{4}\right)^k$. Обозначим объединение этих отрезков $F_k$. Тогда $F=\bigcap_{k=1}^{\infty} F_k$. Мера этого множества равна 0 , так как $\left(\frac{2}{3}\right)^k \rightarrow 0, k \rightarrow \infty$.

    Возникает вопрос: а не выкинем ли мы вообще все точки при таком построении множества? Очевидно нет, так как точно останутся точки 0 и 1. Все точки в множестве мы можем пронумеровать в виде двоичных чисел. На $k$ - шаге мы даем каждому отрезку номер - двоичное число длины $k$.

    Теперь мы можем построить канторову лестницу. Подробное описание построения есть в книге «Действительный анализ в задачах». Функция также строится итеративно. На каждой итерации мы делим отрезок, полученный на предыдущей итерации на три части. Приведем пример, как строится первая итерация для объяснения алгоритма построения. На первой итерации отрезок $[0,1]$ делится на 3 части: $\left[0, \frac{1}{3}\right],\left[\frac{1}{3}, \frac{2}{3}\right],\left[\frac{2}{3}, 1\right]$. Первому отрезку сопоставляем значение 0 (начала разбиваемого отрезка по $y$), второму отрезку сопоставляем значение $\frac{1}{2}$ — середина, третьему отрезку сопоставляем значение 1. Далее мы берем отрезки, полученные на предыдущей итерации и работаем с ними аналогичным образом, «сужая» отрезки и сопоставляемые им множества. Точки, задающие границы отрезки мы соединяем прямыми.

    Таким образом, мы получаем последовательность функций $F_n(x) \rightarrow F(x)$. При этом, $F(1)-F(0)=$ 1, но $F^{\prime}(x)=0$ почти всюду. Получаем, что канторова лестница непрерывна, ограниченной вариации, но не абсолютно непрерывна.
\end{example}


\subsection{Пространство $L_\infty$}
\begin{definition}
    Будем называть \textit{существенным супремумом} для измеримой функции $f$ на пространстве с мерой $(X, \mathfrak{M}, \mu)$ следующее:
    \[\esup\limits_{x \in X} |f(x)| = \inf\{L > 0 \ | \ \mu(|f| > L) = 0\}.\]
\end{definition} 

\begin{definition}
    Пусть $(X, \mathfrak{M}, \mu)$ — пространство с мерой.
    Будем говорить, что $f: X \to \overline{\R}$ \textit{существенно ограничена}, если
    \[\esup\limits_{x \in X} |f(x)| < +\infty.\]
\end{definition} 

\begin{definition}
    Определим $L_\infty$ как ЛНП классов эквивалентности существенно ограниченных на $X$ функций с нормой $\esup$, то есть
    \[\|f\|_\infty = \esup\limits_X |f|.\]
\end{definition} 

\begin{theorem}[Обобщение неравенства Гёльдера]
    Пусть $f\in L_\infty(X)$ и $g\in L_1(X)$. Тогда $fg\in L_1(X)$ и
    \[\|fg\|_1 \le \|f\|_\infty \|g\|_1.\]
\end{theorem}
\begin{proof}
    По определению существенного супремума существует множество $E\subset X$ такое, что
    \[\mu(E)=0\]
    и для всех $x\in X\setminus E$ выполнено
    \[|f(x)|\leqslant \|f\|_\infty.\]

    Тогда почти всюду
    \[|f(x)g(x)|\leqslant \|f\|_\infty |g(x)|.\]
    Интегрируя это неравенство, получаем
    \[\int_X |f g|\,d\mu
    \leqslant
    \int_X \|f\|_\infty |g|\,d\mu
    =
    \|f\|_\infty \int_X |g|\,d\mu.\]
    Следовательно,
    \[\|fg\|_1\leqslant \|f\|_\infty\|g\|_1.\]
\end{proof}

Это предельный случай неравенства Гёльдера при сопряжённых показателях
\[p=\infty,\qquad q=1.\]

Аналогично, если $f\in L_1(X)$ и $g\in L_\infty(X)$, то
\[\|fg\|_1\leqslant \|f\|_1\|g\|_\infty.\]

Также верен вариант для двух существенно ограниченных функций:
если $f,g\in L_\infty(X)$, то
\[\|fg\|_\infty\leqslant \|f\|_\infty\|g\|_\infty.\]

Действительно, почти всюду
\[
|f(x)|\leqslant \|f\|_\infty,
\qquad
|g(x)|\leqslant \|g\|_\infty,
\]
поэтому почти всюду
\[
|f(x)g(x)|
\leqslant
\|f\|_\infty\|g\|_\infty.
\]
Отсюда
\[
\esup_{x\in X}|f(x)g(x)|
\leqslant
\|f\|_\infty\|g\|_\infty.
\]

\begin{theorem}[Обобщение неравенства Минковского]
    Пусть $f,g\in L_\infty(X)$. Тогда $f+g\in L_\infty(X)$ и
    \[
    \|f+g\|_\infty
    \leqslant
    \|f\|_\infty+\|g\|_\infty.
    \]
\end{theorem}
\begin{proof}
    По определению существенного супремума существуют множества $E_f,E_g\subset X$ такие, что
    \[\mu(E_f)=\mu(E_g)=0,\]
    причём при $x\notin E_f$
    \[|f(x)|\leqslant \|f\|_\infty,\]
    а при $x\notin E_g$
    \[|g(x)|\leqslant \|g\|_\infty.\]
    Положим
    \[E:=E_f\cup E_g.\]
    Тогда
    \[\mu(E)=0.\]
    Для всех $x\in X\setminus E$ имеем
    \[
    |f(x)+g(x)|
    \leqslant
    |f(x)|+|g(x)|
    \leqslant
    \|f\|_\infty+\|g\|_\infty.
    \]
    Следовательно,
    \[
    \operatorname{ess\,sup}_{x\in X}|f(x)+g(x)|
    \leqslant
    \|f\|_\infty+\|g\|_\infty.
    \]
    То есть
    \[
    \|f+g\|_\infty
    \leqslant
    \|f\|_\infty+\|g\|_\infty.
    \]
\end{proof}

\begin{theorem}
Пространство $L_\infty(X)$ полно.
\end{theorem}
\begin{proof}
    Пусть $\{f_n\}_{n=1}^\infty$ — фундаментальная последовательность в $L_\infty(X)$. Это означает, что
    \[
    \forall \varepsilon>0\ \exists N\in\mathbb N:
    \forall n,m\geqslant N
    \quad
    \|f_n-f_m\|_\infty<\varepsilon.
    \]
    Нужно доказать, что существует $f\in L_\infty(X)$ такая, что
    \[
    \|f_n-f\|_\infty\to 0,
    \qquad n\to\infty.
    \]

    Зафиксируем $k\in\mathbb N$. По фундаментальности существует $N_k\in\mathbb N$ такое, что для всех $n,m\ge N_k$
    \[\|f_n-f_m\|_\infty<\frac1k.\]
    Это значит, что для каждой пары $n,m\ge N_k$ существует множество $E_{n,m}^{(k)}\subset X$ меры ноль такое, что для всех $x\notin E_{n,m}^{(k)}$
    \[|f_n(x)-f_m(x)|\leqslant \frac1k.\]
    Положим
    \[E_k:=\bigcup_{n,m\geqslant N_k} E_{n,m}^{(k)}.\]
    Так как это счётное объединение множеств меры ноль, то
    \[\mu(E_k)=0.\]
    Теперь положим
    \[E:=\bigcup_{k=1}^{\infty}E_k.\]
    Тогда снова
    \[\mu(E)=0.\]

    Покажем, что для каждого $x\in X\setminus E$ числовая последовательность $\{f_n(x)\}$ фундаментальна. Действительно, для любого $k\in\mathbb N$ и любых $n,m\geqslant N_k$ имеем
    \[|f_n(x)-f_m(x)|\leqslant \frac1k.\]
    Следовательно, $\{f_n(x)\}$ фундаментальна в $\mathbb R$ или $\mathbb C$. Значит, она сходится.

    Определим функцию $f$ следующим образом:
    \[
    f(x):=\lim_{n\to\infty} f_n(x),
    \qquad x\in X\setminus E.
    \]
    На множестве $E$ доопределим $f$ произвольно, например положим
    \[
    f(x):=0,
    \qquad x\in E.
    \]
    Так как $f$ почти всюду является поточечным пределом измеримых функций, то $f$ измерима.

    Теперь докажем, что $f\in L_\infty(X)$. Так как $\{f_n\}$ фундаментальна в $L_\infty$, существует $N\in\mathbb N$ такое, что для всех $n\geqslant N$
    \[\|f_n-f_N\|_\infty<1.\]
    Значит, вне некоторого множества меры ноль выполнено
    \[|f_n(x)-f_N(x)|\leqslant 1.\]
    Переходя к пределу при $n\to\infty$, получаем почти всюду
    \[|f(x)-f_N(x)|\leqslant 1.\]
    Тогда почти всюду
    \[
    |f(x)|
    \leqslant
    |f_N(x)|+1.
    \]
    Следовательно,
    \[
    \|f\|_\infty
    \leqslant
    \|f_N\|_\infty+1
    <+\infty.
    \]
    Значит,
    \[f\in L_\infty(X).\]

    Осталось доказать, что $f_n\to f$ по норме $L_\infty$. Пусть $\varepsilon>0$. Так как $\{f_n\}$ фундаментальна, существует $N\in\mathbb N$ такое, что для всех $n,m\geqslant N$
    \[\|f_n-f_m\|_\infty<\varepsilon.\]
    Значит, вне некоторого множества меры ноль для всех $n,m\geqslant N$ выполнено
    \[|f_n(x)-f_m(x)|\leqslant \varepsilon.\]
    Фиксируем $n\geqslant N$ и переходим к пределу при $m\to\infty$. Тогда почти всюду
    \[|f_n(x)-f(x)|\leqslant \varepsilon.\]
    Следовательно,
    \[\|f_n-f\|_\infty\leqslant \varepsilon\]
    для всех $n\geqslant N$. Поэтому
    \[
    \|f_n-f\|_\infty\to 0,
    \qquad n\to\infty.
    \]

    Итак, любая фундаментальная последовательность в $L_\infty(X)$ сходится по норме $L_\infty$ к элементу из $L_\infty(X)$. Следовательно, пространство $L_\infty(X)$ полно.
\end{proof}